% TFM-ML-cap01-introduccion.tex
% Capítulo 01 — Introducción
% Autor: Mario Lourido Regueira
% Scope: v5 — Corredor CHUF–Ferrol, horizonte variable
% Última revisión: 2026-04-30
% Estado: REDACTADO — pendiente revisión directores


% ============================================================
% 1.1 PRESENTACIÓN DEL PROYECTO
% ============================================================

\subsection{Presentación del proyecto}
\label{sec:presentacion}

Este trabajo de fin de máster propone el concepto de servicio de movilidad nocturna 
para el transporte nocturno de trabajadores de turno en el corredor del Complexo 
Hospitalario Universitario de Ferrol (\textsc{chuf}, A~Cabana) hacia las zonas 
residenciales de la ciudad, operando en la franja horaria comprendida entre las 22:00
 y las 06:00~h; el horizonte temporal se determina en función de la tecnología seleccionada 
 en la fase de diseño.

El trabajo se desarrolla desde la perspectiva del diseño industrial aplicado
al diseño de servicio. El artefacto resultante no es un vehículo ni una
aplicación, sino un sistema de servicio: sus recursos finales son un
\textit{service blueprint} (mapa de servicio), un \textit{journey map}
(mapa de experiencia de usuario), los principios de diseño del espacio
interior del servicio y un modelo de replicabilidad a corredores europeos
análogos.

El usuario principal del servicio es el personal sanitario de turno nocturno
del \textsc{chuf} (enfermería, auxiliares de enfermería y celadores) un
colectivo mayoritariamente femenino que trabaja en una instalación periférica
sin cobertura de transporte público regular en horario nocturno. Este perfil
 es representativo de decenas de ciudades europeas de
tamaño medio con hospital periférico y con deficiencias en el transporte nocturno.
Ferrol actúa aquí como laboratorio de diseño; el modelo de servicio generado
es el podría ser transferible a otras ciudades homólogas.

Lo que este trabajo \emph{no} es: no es el diseño de ingeniería del
vehículo, no es el desarrollo de una aplicación funcional, no es un plan
de negocio cerrado ni la implementación operativa del servicio. Es la
justificación estratégica y el concepto de diseño que habilita cualquiera
de esos desarrollos posteriores.


% ============================================================
% 1.2 MOTIVACIÓN Y JUSTIFICACIÓN
% ============================================================

\subsection{Motivación y justificación}
\label{sec:motivacion}

\begin{claimS}
El trabajo a turnos afecta a una fracción significativa de la fuerza
laboral europea: el 21\,\% de los trabajadores de la UE-28 trabaja en
turno, y el 48,8\,\% de ese colectivo lo hace en turnos rotativos o
alternos que incluyen la noche \cite{SanchezSellero2025}. Entre los
colectivos que expresan con mayor frecuencia la necesidad de movilidad
nocturna fiable se encuentran los trabajadores sanitarios, el personal
de limpieza y mantenimiento y los empleados de hostelería: sectores que
operan en instalaciones periféricas durante la franja horaria en que el
transporte colectivo ha cesado \cite{URBACT2025Gender}. Para todos
ellos, la ausencia de servicio nocturno no es una carencia puntual sino
un coste estructural que se repite cada jornada.
\end{claimS}

\begin{claimS}
El sector sanitario es uno de los principales generadores de movilidad
laboral nocturna estructurada en Europa, caracterizado por turnos
predecibles, horario regulado y origen--destino constantes
\parencite{Eurofound2020Sectors}. En el cuarto trimestre de 2019, 14,7 millones de personas
estaban empleadas en ocupaciones sanitarias en la Unión Europea, de las cuales
el 78\% eran mujeres \parencite{Eurostat2020Health}. En España, el
11,4\% de los ocupados trabajó en jornada nocturna en 2024, con una brecha de
género significativa: el 13,7\% de los hombres frente al 8,8\% de las
mujeres \parencite{INE2025EPA}.
\end{claimS}

La combinación de estos dos factores (vacío nocturno en el transporte
público y fuerza laboral mayoritariamente femenina trabajando en turno de
noche) define un problema de diseño que este
trabajo trata como eje transversal. 

\begin{claimS}
La investigación empírica documenta una brecha cuantificable: las mujeres
tienen una probabilidad un 10\% superior a la de los hombres de sentirse
inseguras en el metro, y un 6\% superior en autobús
\parencite{AitBihiOuali2020}. El 75\% de las mujeres encuestadas modifica
su ruta al desplazarse de noche respecto al día, y solo el 2\% declara
sentirse completamente segura viajando en horario nocturno
\parencite{URBACT2025Gender}.
\end{claimS}

\begin{claimS}
El patrón no es exclusivo de grandes metrópolis. En Braga (Portugal,
aproximadamente 200.000 habitantes), el servicio de autobús urbano
finaliza a medianoche; la propia operadora estudia adaptar los horarios
para los trabajadores del hospital local que residen en zonas suburbanas
sin cobertura posterior \cite{dAntonio2024}. La escala del problema en
ciudades medias es documentable y replicable: no se trata de un déficit
excepcional sino de un patrón estructural que se repite en corredores
hospitalarios periféricos en distintos países europeos.
Desde el lado tecnológico, los proyectos europeos financiados con fondos
H2020 han demostrado la viabilidad operativa de shuttles autónomos
SAE~L4 en entornos urbanos reales. El proyecto \textsc{avenue} (2018--2022)
desplegó servicios regulares en 10 emplazamientos europeos, de los cuales
cuatro mantienen operación comercial tras el fin de la financiación
\parencite{AVENUE2023}. El proyecto \textsc{show} (2020--2024) consolidó
más de 72 vehículos automatizados en 22 ciudades piloto
\parencite{SHOW2024Legacy}. Ambos proyectos fueron diseñados explícitamente
para operar en zonas de demanda baja a media \parencite{AVENUE2023}:
exactamente el perfil de una ciudad como Ferrol.
\end{claimS}

El argumento central de este trabajo es que el servicio de movilidad nocturna 
no es una solución pensada para metrópolis: es estructuralmente adecuado para corredores
de demanda predecible y baja, donde el transporte nocturno convencional
resulta inviable económicamente. Ferrol, con su corredor \textsc{chuf}--zonas
residenciales, es un caso representativo y replicable de ese nicho.


% ============================================================
% 1.3 OBJETIVOS
% ============================================================

\subsection{Objetivos}
\label{sec:objetivos}

El \textbf{objetivo general} de este trabajo es diseñar un concepto de servicio de 
movilidad nocturna para trabajadores sanitarios en el corredor CHUF–Ferrol, 
fundamentado en evidencia empírica propia, y extraer de él un modelo de servicio 
transferible a ciudades europeas con perfil análogo.

Los \textbf{objetivos específicos} son:

\begin{enumerate}[label=\arabic*., leftmargin=*, noitemsep]
  \item Analizar el vacío de transporte nocturno para trabajadores esenciales
        en ciudades europeas de tamaño medio y documentar su impacto en la
        experiencia de movilidad y la carga de fatiga del trabajador nocturno.
  \item Caracterizar el perfil del usuario afectado y documentar empíricamente
        sus barreras de movilidad nocturna mediante investigación primaria
        (encuesta propia, n=167, provincia de A~Coruña), con literatura
        europea como marco de contextualización.
  \item Diseñar el concepto de servicio completo: \textit{service blueprint},
        \textit{journey map} y principios de diseño del espacio interior del
        servicio.
  \item Argumentar la transferibilidad del modelo de servicio a corredores
        europeos análogos: ciudades de 50.000--150.000 habitantes con
        hospital periférico y franja sin cobertura de transporte nocturno.
\end{enumerate}


% ============================================================
% 1.4 ALCANCE Y LIMITACIONES
% ============================================================

\subsection{Alcance y limitaciones}
\label{sec:alcance}

El trabajo se circunscribe al \textbf{diseño de servicio y a la propuesta
conceptual}. El ámbito geográfico es doble: el corredor \textsc{chuf}--Ferrol
como caso de estudio, y las ciudades europeas de tamaño medio como ámbito
del modelo replicable.

El horizonte temporal se determina en función de la tecnología seleccionada en la fase de diseño, fundamentándose en los datos actuales de operación de servicios nocturnos y en las proyecciones de madurez de las soluciones candidatas.

\textbf{Quedan fuera del alcance:}
\begin{itemize}
  \item El diseño de ingeniería del vehículo (carrocería, tren de potencia,
        sistemas de percepción).
  \item El desarrollo de una aplicación funcional o interfaz codificada.
  \item El análisis financiero de detalle (cuenta de resultados, flujo
        de caja).
  \item La implementación operativa del servicio.
  \item Las verticales de transporte nocturno urbano convencional, que se utilizan únicamente como marco
        contextual.
\end{itemize}

Las \textbf{limitaciones} del trabajo son de naturaleza metodológica y de
alcance. Primero, el trabajo se apoya en un caso de estudio único: las
conclusiones sobre replicabilidad se argumentan por analogía estructural
con corredores europeos comparables, pero no se validan empíricamente sobre
múltiples casos. Por último, el horizonte temporal se determina en función de la tecnología seleccionada en la fase de diseño, lo que implica que el concepto de servicio es prospectivo: la tecnología seleccionada puede no operar de manera madura comercialmente en Europa en la fecha de redacción, y las proyecciones de limitaciones regulatorias y tecnológicas están sujetas a incertidumbre inherente.


% ============================================================
% 1.5 METODOLOGÍA
% ============================================================

\subsection{Metodología}
\label{sec:metodologia}

El enfoque metodológico es el \textbf{diseño de servicio aplicado desde la
ingeniería en diseño industrial}. El trabajo combina investigación de
escritorio (\textit{desk research}) con análisis de datos empíricos de
terceros, investigación de campo (Entrevistas y encuestas), siguiendo una filosofía de diseño centrado en el usuario en la
que el usuario informa y el diseñador decide \parencite{MartinezRodriguez2024}.

Las \textbf{herramientas de aplicación confirmada} son:

\begin{itemize}[noitemsep]
  \item Análisis \textsc{pestel} de la movilidad nocturna en Europa.
  \item \textit{Stakeholder mapping} por matriz Poder--Interés: actores del ecosistema de movilidad nocturna en el contexto del corredor Ferrol.
  \item Despliegue de la función de calidad (\textsc{qfd}): traducción
        de necesidades del usuario en requisitos del servicio y del
        espacio interior.
  \item \textit{Benchmarking} de servicios de movilidad nocturna en corredores análogos.
  \item \textit{Service blueprint}: mapa completo del servicio en el
        corredor \textsc{chuf}--Ferrol (acciones del usuario,
        \textit{frontstage}, \textit{backstage}, procesos de soporte).
  \item \textit{Journey map}: experiencia del usuario en todas las fases
        del viaje (pre-viaje, espera, embarque, trayecto, llegada).
  \item \textit{Business Model Canvas}: modelo de negocio conceptual
        del servicio.
  \item Priorización \textsc{moscow} de requisitos del servicio.
  \item Matriz de afirmaciones--evidencia (\textit{Supported} /
        \textit{Contested} / \textit{Undetermined}) con nivel de
        confianza (Alta / Media / Baja).
  \item Análisis \textsc{dafo} del sistema propuesto.
\end{itemize}

El tratamiento de los datos sigue un protocolo explícito de verificación de
fuentes: cada afirmación factual de la memoria lleva asociada su fuente
externa verificable, con autor, año y \textsc{doi} o \textsc{url}.

La investigación empírica se procesa mediante una cadena de síntesis
inductiva en cuatro eslabones: el \textbf{dato} cuantitativo o cualitativo
verificado, la \textbf{observación} que interpreta el patrón subyacente,
el \textbf{\textit{insight}} que formula la necesidad latente del usuario,
y el \textbf{criterio de diseño} que traduce esa necesidad en un requisito
del servicio. Ningún criterio de diseño se formula sin que haya un dato
verificado que lo origine.

% ============================================================
% 1.6 ESTRUCTURA DE LA MEMORIA
% ============================================================

\subsection{Estructura de la memoria}
\label{sec:estructura}

La memoria se organiza en diez capítulos.

El \textbf{capítulo 2} establece el marco contextual: el análisis \textsc{pestel} de la movilidad nocturna en Europa, el estado de la tecnología de vehículos autónomos, el caso de estudio de Ferrol y el mapa de \textit{stakeholders} del ecosistema.

El \textbf{capítulo 3} contiene la investigación de usuario: la encuesta primaria propia (n=167), el análisis cuantitativo, el análisis cualitativo de respuestas abiertas y la perspectiva de género en el riesgo nocturno.

El \textbf{capítulo 4} recoge el análisis de mercado y tendencias: el \textit{benchmarking} de pilotos europeos de movilidad nocturna y el análisis de tendencias del sector.

El \textbf{capítulo 5} sintetiza los cinco \textit{insights} de diseño derivados de la investigación, estructurados como tensión usuario--sistema, dato, observación y criterio de diseño.

El \textbf{capítulo 6} desarrolla las estrategias de diseño: la traducción de necesidades en requisitos mediante \textsc{qfd} y \textsc{moscow}, los valores del servicio y el grado de innovación.

El \textbf{capítulo 7} presenta el \textsc{adn} del concepto: la arquitectura del servicio, la evaluación de candidatos tecnológicos, el \textit{service blueprint}, el \textit{journey map}, los \textit{touchpoints} y el modelo de negocio conceptual.

El \textbf{capítulo 8} recoge la evaluación del proyecto: el pliego de condiciones, el análisis \textsc{dafo}, la matriz de afirmaciones--evidencia, el valor diferencial y el análisis de viabilidad.

El \textbf{capítulo 9} presenta las conclusiones: síntesis de hallazgos, limitaciones metodológicas y líneas futuras de investigación.

El \textbf{capítulo 10} detalla el presupuesto del proyecto como ejercicio de cuantificación del coste real del trabajo de diseño.